\section{Introduction}

Throughout, \(\N=\{0,1,2,\ldots\}\), \(\log\) denotes the natural
logarithm, and
\[
                         \theta=\frac{\log 3}{\log 2}.
\]
For an irrational real number \(\xi\), its irrationality exponent
\(\mu(\xi)\) is the supremum of the real numbers \(\nu\) for which
\[
                 0<\left|\xi-\frac pq\right|<q^{-\nu}
\]
has infinitely many reduced solutions \(p/q\) with \(q>0\).
Unique factorization shows that \(\theta\) is irrational: a rational
identity \(\theta=p/q\), with \(p,q>0\), would imply \(3^q=2^p\).
Dirichlet's approximation theorem therefore gives \(\mu(\theta)\geq2\).
We prove the following upper bound.

\begin{theorem}\label{thm:main}
For every \(\varepsilon>0\), there exists an integer
\(Q_\varepsilon\geq2\) such that, for every \(p\in\Z\) and every
integer \(q\geq Q_\varepsilon\),
\begin{equation}\label{eq:main}
 \left|\frac{\log 3}{\log 2}-\frac pq\right|
                       \geq q^{-3-\varepsilon}.
\end{equation}
Consequently,
\[
                         2\leq\mu(\theta)\leq3.
\]
\end{theorem}

Rational approximation to \(\theta\) is equivalently a homogeneous
linear-form problem:
\[
 \left|\theta-\frac pq\right|
       =\frac{|q\log 3-p\log 2|}{q\log 2}.
\]
For comparison, the linear-form estimate of Wu and Wang
\cite{WuWang2014} implies \(\mu(\theta)\leq5.1163051\).
The author's earlier preprint \cite{JhaEarlier} gives the bound
\(5.104104\), by a different construction with symmetrized
polynomials and exact finite certificates. Theorem~\ref{thm:main}
improves this numerical upper bound to three.

The geometric and analytic framework is inspired by OpenAI's proof
of the exponent-two theorem for \(\pi\) \cite{OAIpi}, in particular
its use of separated weights, interpolation, and determinants of
common Taylor coefficients. Liu, Jiang, and Zhang \cite{LJZ2026}
establish exponent two for every nonzero logarithm of a positive
rational number. The author's rational-logarithm and
compatible-embedding results are treated separately in
\cite{JhaLogarithms}. A bound for each individual logarithm does not
by itself give the corresponding bound for their quotient.
Here both base variables are retained, and primitive rational
characters provide the curve rigidity needed for two-base interpolation.
The final determinant estimate counts auxiliary and Taylor indices
together.

\subsection{Outline and order of choices}

It suffices to exclude reduced approximations
\begin{equation}\label{eq:exceptional}
                         |\theta-p/q|\leq q^{-\nu}
\end{equation}
with unbounded denominators for every rational \(\nu>3\).
Assuming such approximations exist, we first fix rational seed
parameters, then a finite auxiliary dimension and base weights.
Only after these choices do we select finitely many exceptional
fractions with sufficiently separated denominators.
All these data remain fixed while the interpolation degree \(H\)
tends to infinity.

Sections~\ref{sec:local}--\ref{sec:geometry} establish the
interpolation theorem from a local multiplicity comparison and a
global curve inequality. Section~\ref{sec:arithmetic} extracts a
nonzero rational determinant and gives its arithmetic lower bound.
Sections~\ref{sec:analysis}--\ref{sec:direct-packing} obtain an
incompatible upper bound by translating the rows to common analytic
functions and counting distinct Taylor indices.
Section~\ref{sec:parameters} chooses the parameters and completes
the proof. Section~\ref{sec:verification} describes the formal
verification and its reproducible source package.

The conclusion is an upper bound, not an equality with three.
The argument does not provide an effective numerical value of
\(Q_\varepsilon\).
